\section{Conclusion}
\hspace*{1cm} HeatItOn is a utility in the city of Heatington that supplies heat to 
approximately 1,600 buildings through a centralized district heating network. Prior to 
this project, heat production was planned manually, making it difficult to consistently 
select the most efficient combination of production units. The goal was to replace this 
manual process with an automated optimization system capable of generating heat 
production schedules that consistently meet heat demand while minimizing operational 
costs.

\hspace*{1cm} The system was built around five core modules: the Optimizer, Data 
Visualization, Asset Manager, Source Data Manager, and Result Data Manager. The 
architecture followed a clear separation between the Frontend, Backend, and Database 
layers.

\hspace*{1cm} Both required scenarios were successfully implemented. Scenario 1 covered 
a heat-only network, where the system prioritized cheaper units and handled maintenance 
scheduling. Scenario 2 introduced electricity production, where the optimizer calculated 
net production costs on an hourly basis to determine when to sell or consume electricity. 
A custom scenario was also implemented, giving users control over which assets are 
included in the optimization.

\hspace*{1cm} The final product successfully addressed the original problem statement. 
The system is able to automatically generate optimized heat production schedules that 
meet the heat demand of Heatington while minimizing operational costs, replacing the 
previous manual planning process entirely. Operators can now manage assets, import 
source data, run the optimizer, and view the results through a single integrated 
application.

\hspace*{1cm} Looking ahead, the system could be further extended in several ways. 
Integrating live electricity prices from Energinet.dk would allow the optimizer to react 
to real-time market data. The ability to enable or disable individual units directly from 
the user interface would give operators greater flexibility during optimization. In future 
development, the team would also focus on improving test coverage, refining the user 
interface based on user feedback, and expanding the system to support additional heating 
scenarios.
